\documentclass{article}
\usepackage[utf8]{inputenc}
\usepackage[T1]{fontenc}
\usepackage{geometry}
\geometry{
    a4paper,
    left=2.5cm,
    right=2.5cm,
    top=2.5cm,
    bottom=2.5cm
}
\usepackage{graphicx}
\usepackage{float}
\usepackage{hyperref}
\usepackage[english]{babel}
\usepackage{amsmath}
\usepackage{amssymb}
\usepackage{array}
\usepackage{booktabs}
\usepackage{caption}
\usepackage{subcaption}
\usepackage{enumitem}
\usepackage{fancyhdr}
\usepackage{titlesec}
\usepackage{url}
\usepackage{textcomp}
\usepackage{xcolor}
\usepackage{listings}
\usepackage{multirow}
\usepackage{longtable}
\usepackage{tabularx}
\usepackage{times}

% Set up page style
\pagestyle{fancy}
\fancyhf{}
\rfoot{\thepage}
\renewcommand{\headrulewidth}{0pt}

% Title formatting
\titleformat{\chapter}[display]{\normalfont\bfseries}{\chaptertitle\thechapter}{0pt}{\Huge}
\titlespacing*{\chapter}{0pt}{50pt}{40pt}

% Document info
\title{Template Document}
\author{Academic Support Agent}
\date{\today}

\begin{document}
\maketitle

\begin{frame}
	\titlepage
\end{frame}

\begin{frame}
	
\end{frame}

\section{A section}
\subsection{With a subsection}

\begin{frame}
	\frametitle{There Is No Largest Prime Number}
	\framesubtitle{The proof uses \textit{reductio ad absurdum}.}
	\begin{theorem}
		There is no largest prime number.
	\end{theorem}
	\begin{proof}
		\begin{enumerate}
			\item<1-| alert@1> Suppose $p$ were the largest prime number.
			\item<2-> Let $q$ be the product of the first $p$ numbers.
			\item<3-> Then $q+1$ is not divisible by any of them.
			\item<1-> But $q + 1$ is greater than $1$, thus divisible by some prime
				number not in the first $p$ numbers.\qedhere
		\end{enumerate}
	\end{proof}
\end{frame}

\begin{frame}
	\frametitle{Block colors}
	\begin{block}{A block}
		With text
	\end{block}
	\begin{alertblock}{An alert block}
		With text
	\end{alertblock}
	\begin{exampleblock}{An example block}
		\begin{itemize}
			\item An item
			\item And another one
		\end{itemize}
	\end{exampleblock}
\end{frame}
\end{document}